% Chapter 1, Topic from _Linear Algebra_ Jim Hefferon
%  http://joshua.smcvt.edu/linearalgebra
%  2001-Jun-09
\topic{Accuracy of Computations}
\index{accuracy! of Gauss's Method|(}
\index{Gauss's Method! accuracy|(}
Gauss's Method lends itself to computerization.
This C~language\index{C language} 
code illustrates.
It operates on an $\nbyn{n}$ matrix named~\lstinline[style=inline]!a!.
\begin{lstlisting}
for(row=1; row<=n-1; row++){
  for(row_below=row+1; row_below<=n; row_below++){
    multiplier = a[row_below,row]/a[row,row];
    for(col=row; col<=n; col++){
      a[row_below,col] -= multiplier*a[row,col];
    }
  }
}
\end{lstlisting}
The 
\lstinline[style=inline]!for(row=1; row<=n-1; row++){ .. }!
construct in the top line
initializes \lstinline[style=inline]!row! to~$1$ and then iterates while
\lstinline[style=inline]!row! is less than or equal to $n-1$, each time through
incrementing \lstinline[style=inline]!row! by one with the 
\lstinline[style=inline]!++! operation.
The other non-obvious language 
construct is the \lstinline[style=inline]!-=! in the innermost
loop, which produces the same result as
\lstinline[style=inline]!a[row_below,col] = -1*multiplier*a[row,col] + a[row_below,col]!.

That code is a good first take but
it is naive.
For one thing, it assumes that 
the entry in
\lstinline[style=inline]!row,row! is nonzero.
So the code needs to be extended to cover
the case where finding a zero there leads to a row swap or to the
conclusion that the matrix is singular.

We could add \lstinline[style=inline]!if! statements to cover those cases 
but we will instead consider another way in which this code is naive.
It is prone to pitfalls arising from the computer's reliance on 
floating point arithmetic.

A full discussion of floating point would be a detour.
But in short, computer numbers are in a form of scientific notation,
$m\times 10^n$\!. 
Typically $m$ has at most sixteen digits and~$n$ is bounded
between about plus and minus one thousand.
That's quite a range but nonetheless trouble can arise.

Consider this system.
The numbers in the second line have seventeen digits.
\newlength{\jhwidenumberlen}
\settowidth{\jhwidenumberlen}{$3.000\,000\,000\,000\,000\,1$}
\begin{equation*}
   \begin{linsys}{2}
                   x &+ &2y &= &\makebox[\jhwidenumberlen][l]{$3$}  \\
     1.000\,000\,000\,000\,000\,1x &+ &2y &= &3.000\,000\,000\,000\,000\,1
   \end{linsys}
\end{equation*}
By eye we spot the solution $x=1$, $y=1$.
A computer
using floating point will represent the second
equation as 
$1.000\,000\,000\,000\,000x+2y=3.000\,000\,000\,000\,000$, losing the
digits in the last place.
It will report that the system is 
singular.

That example is unwieldy.
Below we illustrate with examples using only a few digits;~tweaking to apply to 
more digits is easy.

For some intuition into the prior example,
consider the next system.
\begin{equation*}
   \begin{linsys}{2}
         x &+ &2y &= &3\hfill\hbox{}  \\
     1.01x &+ &2y &= &3.01
   \end{linsys}
   \tag{$*$}
\end{equation*}
We have seen that we must handle a singular system as
a separate case.
But this system is nearly singular in that
the two lines are nearly the same, and also requires special care.
Below on the left, the equation $x+2y=3$ 
is graphed as a dashed line, while the other equation is solid.
Looking very close, we can only barely tell the two apart.
\begin{center}
  % \includegraphics{gr/mp/ch1.33}
  \includegraphics{gr/asy/gr000.pdf}
  \hspace*{3em}
  \includegraphics{gr/asy/gr001.pdf}
\end{center}
The tiny angle between the lines
means that a small change in an
equation causes a large change in the solution.
For instance, switching $3.01$ to $3.03$ 
moves the intersection from $(1,1)$ to $(3,0)$.
The solution
changes radically depending on the final digit, which explains why
a computer with a limited number of decimal places has trouble.
% sage: eqns = [x+2*y==3, 5.01*x+2*y==3.01]
% sage: solve(eqns, x, y)
% [[x == (1/401), y == (601/401)]]
% sage: eqns = [x+2*y==3, 5*x+2*y==3]
% sage: solve(eqns, x, y)
% [[x == 0, y == (3/2)]]
(This sensitivity disappears if the angle is not small.
If ($*$)'s second equation were $5.01x + 2y = 3.01$ then dropping the 
$.01$'s would change the solution hardly at all, 
from $(1/401,601/401)$ to $(0,3/2)$.)
A problem that is very sensitive to changes in
the input values is \definend{ill-conditioned}.\index{ill-conditioned problem}

That example shows one way in which a system can give
a computer trouble.
Its advantage is that the picture of nearly-equal lines is memorable.
Its disadvantage is that in solving
large systems we usually do not understand much about the geometry.

There are other ways that the computer's results may be
unreliable.
Consider this system (from \cite{Hamming}).
The graph shows that it is far from singular.
\begin{equation*}
  \begin{linsys}{2}
     0.001x  &+  &y  &=  &1  \\
          x  &-  &y  &=  &0
  \end{linsys}
  \hspace*{3em}
  \begin{tabular}{c}\includegraphics{gr/asy/gr002.pdf}\end{tabular}
\tag*{($**$)}\end{equation*}
The second equation
gives $x=y$, so $x=y=1/1.001$ and 
thus we know that the correct answer is that
both variables have values that are just less than $1$.

A computer using three digits represents the system internally in this way.
\begin{equation*}
  \begin{linsys}{2}
    (1.00\times 10^{-3})\cdot x  &+  &(1.00\times 10^{0})\cdot y  &=  &1.00\times 10^{0}  \\
    (1.00\times 10^{0})\cdot x   &-  &(1.00\times 10^{0})\cdot y  &=  &0.00\times 10^{0}
  \end{linsys}
\end{equation*}
The Gauss's Method step $-1000\rho_1+\rho_2$ produces 
a second equation $-1001y=-1000$, which rounds to 
$(-1.00\times 10^{3})y=-1.00\times 10^{3}$\!.
The computer decides that $y=1$ 
and with that it concludes from the first equation that $x=0$.
The~$y$ value is close but the~$x$ is bad\Dash
the ratio of the actual answer to the computer's answer is 
infinite.\index{accuracy!rounding error}
In short, another cause of 
unreliable output is using leading entries that are small. 

We can obviously tweak the above to give the
same trouble in higher digits so we need a strategy.
A basic improvement on the earlier naive code 
is to not determine the factor to use for row combinations
by simply taking the entry in the \lstinline[style=inline]!row,row! position,
but rather to also look at the entries in that 
column and below that entry,
and take one that is not too small.
This is \definend{partial pivoting}.%
\index{partial pivoting}\index{pivoting!full} 

For example, to solve the troublesome system~($**$) above
we can start by looking in the first column 
of both equations.
We select the~$1$ in
the second equation as more likely to give good results.
The reduction step $-0.001\rho_2+\rho_1$ gives the first equation 
$1.001y=1$, which the computer represents as 
$(1.00\times 10^{0})y=1.00\times 10^{0}$\!. 
It concludes that $y=1$ and after back-substitution that $x=1$, 
both of which are close to right.  
We can adapt the code from above to do this.
\begin{lstlisting}
for(row=1; row<=n-1; row++){
/* find the largest entry in this column (in row max) */
  max=row;
  for(row_below=row+1; row_below<=n; row_below++){
    if (abs(a[row_below,row]) > abs(a[max,row]));
      max = row_below;
  }
/* swap rows to move that best entry up */
  for(col=row; col<=n; col++){
    temp=a[row,col];
    a[row,col]=a[max,col];
    a[max,col]=temp;
  }
/* proceed as before */
  for(row_below=row+1; row_below<=n; row_below++){
    multiplier=a[row_below,row]/a[row,row];
     for(col=row; col<=n; col++){
       a[row_below,col]-=multiplier*a[row,col];
    }
  }
}
\end{lstlisting}

A full analysis of the best way to implement Gauss's Method 
is beyond our scope; see \cite{Wilkinson65}.
But the method recommended by most experts 
first finds the best entry
among the candidates and then scales it to a number that is less
likely to give trouble.
This is 
\definend{scaled partial pivoting}.\index{scaled partial pivoting}\index{pivoting!partial!scaled}

In addition to returning a result that is likely to be reliable,
most well-done code will return a 
\definend{condition number}.%
\index{condition number}\index{matrix!condition number} 
Roughly, it measures how sensitive the answer is to changes in the input data 
and to rounding errors made while computing the solution.
(The formal definition is that it is the asymptotic worst-case relative change 
   in output compared with a change in input.
% Roughly, it describes the factor by which uncertainties in the input numbers
% could be magnified to become inaccuracies in the results returned;~see 
% \cite{Rice}.
There are different condition numbers for
different problems.)

The lesson is that
just because Gauss's Method always works in theory, and just
because computer code correctly implements that method,
doesn't mean that the answer is reliable.
In practice, always use a package
where experts have worked hard to counter what can go wrong.

\begin{exercises}
  \item 
    Using two decimal place numbers, add $253$ and $2/3$.
    \begin{answer}
      Scientific notation is convenient for expressing 
      the two-place restriction:
      $.25\times 10^{2}+.67\times 10^{0}=.25\times 10^{2}$\!.
      Adding the $2/3$ has no effect.
    \end{answer}
  % \item % Put an example in the body
  %   This intersect-the-lines problem contrasts with the example
  %   discussed above.
  %   \begin{center}
  %     \vcenteredhbox{\includegraphics{gr/mp/ch1.34}}
  %     \qquad
  %     $\displaystyle \begin{linsys}{2}
  %           x &+ &2y &= &3  \\
  %           3x &- &2y &= &1
  %     \end{linsys}$
  %   \end{center}
  %   Illustrate that in this system 
  %   some small change in the numbers will produce only a
  %   small change in the solution by changing the constant in the
  %   bottom equation to $1.008$ and solving.
  %   Compare it to the solution of the unchanged system.
  %   \begin{answer}
  %     The reduction
  %     \begin{equation*}
  %       \grstep{-3\rho_1+\rho_2}
  %       \begin{linsys}{2}
  %         x  &+  &2y   &=  &3  \\
  %            &   &-8y  &=  &-7.992
  %       \end{linsys}
  %     \end{equation*}
  %     gives \( (x,y)=(1.002,0.999) \).
  %     So for this system a small change in the constant produces only a small
  %     change in the solution.
  %   \end{answer}
\item 
    Solve this system (from \cite{Rice}) by rounding at each step to four digits.
    \begin{equation*}
      \begin{linsys}{2}
        0.000\,3x  &+  &1.556y  &=  &1.559 \\
        0.345\,4x  &-  &2.346y  &=  &1.108
      \end{linsys}
    \end{equation*}
    \begin{answer}
          Reduction uses the 
          step $-(0.3454/0.0003)\rho_1+\rho_2$.
          % \begin{equation*}
          %   \grstep{-(.3454/.0003)\rho_1+\rho_2}
          %   \begin{amat}{2}
          %     .0003  &1.556  &1.559  \\
          %     0      &-1794   &-1805
          %   \end{amat}
          % \end{equation*}
          First, $-0.3454/0.0003=-1151.33333333333$ and rounding to 
          four digits gives $-1151=-1.115\times 10^3$\!.
          % sage: factor = 0.3454/0.0003
          % sage: n(factor)
          % 1151.33333333333
          We take that the leading entry in the second row is now zero
          (that's the point of the row reduction step).
          For the second entry, $1.556\cdot(-1151)=-1790.956$, which rounds
          to four digits as $-1791=-1.791\times 10^3$\!.
          Adding $-1791-2.346$ rounds to $-1793=-1.793\times 10^3$\!.
          As to the third entry, $1.559\cdot(-1151)$ rounds to $-1794$,
          and adding rounds to $-1793=-1.793\times 10^3$\!.
          % sage: 1.556*(-1151)
          % -1790.95600000000
          % sage: -1791-2.346
          % -1793.34600000000
          % sage: 1.559*(-1151)
          % -1794.40900000000
          % sage: -1794+1.108
          % -1792.89200000000
          \begin{equation*}
            \begin{amat}{2}
              0.0003  &1.556  &1.559  \\
              0       &-1793  &-1793
            \end{amat}
          \end{equation*}
          The conclusion is that $y$ is $-1793/-1793$, which rounds to 
          $1.000=1.000\times 10^0$\!.
          Substituting gives $1.556\cdot 1.000$, which rounds to 
          $1.556$, and then $1.559-1.556$ rounds to $0.003$.
          Solving $0.0003x=0.003$ gives $0.003/0.0003$, which is 
          $y=10=1.000\times 10^1$\!.
          Checking that by substitution into the original equations shows
          that it is correct.
    \end{answer}
    \item 
    Assume that the machine uses four digit numbers.
    \begin{exparts}
      \partsitem Show that it will compute 
         $(x+y)+z$ as unequal to $x+(y+z)$ 
         when $x=-1.000\times 10^5$, 
         and $y=1.000\times 10^5$,
         and~$z=1.000\times 10^0$.
      \partsitem Compare the computer's versions of $(1/3)x+y=0$
        and $(2/3)x+2y=0$.
        Is twice the first equation the same as the second?
    \end{exparts}
    \begin{answer}
      \begin{exparts}
        \partsitem For the first, $x+y$ is zero and so $(x+y)+z$ is~$1$. 
          For the other, 
          $(y+z)$ is $10\,000+1$ rounded to four digits, $1.000\times 10^5$\!.
          Then $x+(y+z)$ is zero.
        \partsitem The first equation is 
          $0.333\,3\cdot x+1.000\,0\cdot y=0.000\,0$
          while the second is 
          $0.666\,7\cdot x+2.000\,0\cdot y=0.000\,0$.
          Twice the first is 
          $0.666\,6\cdot x+2.000\,0\cdot y=0.000\,0$,
          so they aren't equal.
      \end{exparts}
    \end{answer}
  \item 
    Ill-conditioning is not only dependent on the matrix of
    coefficients.
    This example \cite{Hamming} shows that it can arise from an
    interaction between the left and right sides of the system.
    Let $\varepsilon$ be a small real.
    \begin{equation*}
      \begin{linsys}{3}
        3x  &+  &2y           &+  &z            &=  &6\hfill   \\
        2x  &+  &2\varepsilon y  &+  &2\varepsilon z  
                                          &=  &2+4\varepsilon\hfill \\
         x  &+  &2\varepsilon y  &-  &\varepsilon z   
                                          &=  &1+\varepsilon\hfill
      \end{linsys}
    \end{equation*}
    \begin{exparts}
      \partsitem Solve the system by hand.
        Notice that the $\varepsilon$'s divide out only because there is
        an exact cancellation of the integer parts on the right side
        as well as on the left. 
      \partsitem Solve the system by hand, rounding to two decimal
        places, and with $\varepsilon=0.001$.
%      \partsitem If you have access to a computer package for solving
%        linear systems, see if you can find an $\varepsilon$ small enough
%        to get your package to give incorrect results.
    \end{exparts}
    \begin{answer}
      \begin{exparts}
        \partsitem Here is the calculation.
          \begin{align*}
            &\grstep[-(1/3)\rho_1+\rho_3]{-(2/3)\rho_1+\rho_2}
            \begin{amat}{3}
              3  &2                   &1                   &6               \\
              0  &-(4/3)+2\varepsilon &-(2/3)+2\varepsilon &-2+4\varepsilon \\
              0  &-(2/3)+2\varepsilon &-(1/3)-\varepsilon  &-1+\varepsilon 
            \end{amat}                                                    \\
            &\grstep{-(1/2)\rho_2+\rho_3}
            \begin{amat}{3}
              3  &2                   &1                   &6               \\
              0  &-(4/3)+2\varepsilon &-(2/3)+2\varepsilon &-2+4\varepsilon \\
              0  &\varepsilon         &-2\varepsilon       &-\varepsilon 
            \end{amat}
          \end{align*}
          The third equation is $\varepsilon y-2\varepsilon z=-\varepsilon$.
          Cancel the $\varepsilon$'s and solve to get $y=-1+2z$.
          Substitute into the second equation.
          \begin{align*}
            ((-4/3)+2\varepsilon)(-1+2z)+((-2/3)+2\varepsilon)z
            &=-2+4\varepsilon     \\
            (4/3)-2\varepsilon-(10/3)z+6\varepsilon z
            &=-2+4\varepsilon     \\
            ((-10/3)+6\varepsilon)z
            &= (-10/3)+6\varepsilon
          \end{align*} 
          Thus $z=1$ and so $y=1$.
          Then the first equation says that $x=1$. 
        \partsitem 
          As above, scientific notation is convenient to express the
          restriction on the number of digits.
          With two digits 
          \begin{multline*}
            \begin{amat}{3}
              3.0\times 10^{0}  &2.0\times 10^{0}  &1.0\times 10^{0} 
                 &6.0\times 10^{0}        \\
              1.0\times 10^{0}  &2.0\times 10^{-4} &2.0\times 10^{-4} 
                 &.20\times 10^{1}        \\
              3.0\times 10^{0}  &2.0\times 10^{-4}  &-1.0\times 10^{-4} 
                 &1.0\times 10^{0}        
            \end{amat}                                           \\
            \begin{aligned}
            &\grstep[-(1/3)\rho_1+\rho_3]{-(2/3)\rho_1+\rho_2}
            \begin{amat}{3}
              3.0\times 10^{0}  &2.0\times 10^{0}  &1.0\times 10^{0} 
                 &6.0\times 10^{0}        \\
              0                 &-1.3\times 10^{0} &-6.7\times 10^{-1} 
                 &-2.0\times 10^{0}        \\
              0                 &-6.7\times 10^{-1}  &-3.3\times 10^{-1} 
                 &-1.0\times 10^{0}        
            \end{amat}                                          \\
            &\grstep{-(.67/1.3)\rho_2+\rho_3}
            \begin{amat}{3}
              3.0\times 10^{0}  &2.0\times 10^{0}  &1.0\times 10^{0} 
                 &6.0\times 10^{0}        \\
              0                 &-1.3\times 10^{0} &-6.7\times 10^{-1} 
                 &-2.0\times 10^{0}        \\
              0                 &0                  &1.5\times 10^{-3} 
                 &3.1\times 10^{-3}        
            \end{amat}
            \end{aligned}
          \end{multline*}
          we get $z=2.1$, $y=2.6$, and $x=-0.43$.
          These differ significantly
          from the true values found in the first part.
      \end{exparts}
    \end{answer}
\end{exercises}
\index{accuracy! of Gauss's Method|)}
\index{Gauss's Method! accuracy|)}

\endinput




